\section{与相关工作在实验数据集上的比较}
\label{sec:datasets}
\begin{table}[t]
    \centering
    \begin{tabular}{l|c|c|c|c}
    \toprule
        数据集 & 数据类型 & 词表  & $N$ & 真实世界语言  \\
         \midrule
        Tomita grammars  \cite{tomita1982dynamic}  & 离散 & 2  & 10k & $\times$ \\
        Motzkin grammar \citep{alexander2021exact} & 离散 & 3 & 10k   &$\times$ \\
         Email addresses \citep{radev2008clair} & 离散 & $\leq 256$ & 4k & $\times$  \\
        \midrule
        POWER  \cite{Dua:2019} & 连续 & - & 1659k & $\times$ \\
        GAS \cite{fonollosa2015reservoir} & 连续 & - & 852k &  $\times$ \\
        HEPMASS \cite{baldi2016parameterized} & 连续 & - & 315k & $\times$  \\
        MINIBOONE \cite{roe2005boosted} & 连续 & -  & 29k & $\times$  \\
        BSDS300 \citep{martin2001database}  & 连续 & - & 1000k &  $\times$ \\

        \midrule
        PTB \citep{marcinkiewicz1994building} & 离散 & 10k  & 31k & $\surd$ \\
        WikiText-2  \citep{merity2016pointer} & 离散 & 30k & 73k &  $\surd$ \\

 \bottomrule
    \end{tabular}
    \caption{“数据集”一列表示训练数据集。前三个数据集用于 \citep{miller2021tensor}，接下来的五个用于 \cite{novikov2021tensor}，最后两个用于本文。“$N$”一列表示训练样本的数量。“词表”一列表示词的类型数。“连续”表示连续变量，而“离散”表示离散变量。}
    \label{tab:datasets}
\end{table}



\end{document}


% This document was modified from the file originally made available by
% Pat Langley and Andrea Danyluk for ICML-2K. This version was created
% by Iain Murray in 2018, and modified by Alexandre Bouchard in
% 2019 and 2021 and by Csaba Szepesvari, Gang Niu and Sivan Sabato in 2022.
% Modified again in 2023 by Sivan Sabato and Jonathan Scarlett.
% Previous contributors include Dan Roy, Lise Getoor and Tobias
% Scheffer, which was slightly modified from the 2010 version by
% Thorsten Joachims & Johannes Fuernkranz, slightly modified from the
% 2009 version by Kiri Wagstaff and Sam Roweis's 2008 version, which is
% slightly modified from Prasad Tadepalli's 2007 version which is a
% lightly changed version of the previous year's version by Andrew
% Moore, which was in turn edited from those of Kristian Kersting and
% Codrina Lauth. Alex Smola contributed to the algorithmic style files.
